\chapter{Fork A: the symbolic $(a,c)$-family and the up-to-reparam capstone}
\label{chap:hyperbolic_family}

% Chapter for Gluck/Hyperbolic/Family.lean (aggregator) and its sub-modules
% Gluck/Hyperbolic/Family/{Bicircle, Anchor, Node, Transport, TurningA,
% TurningB, TurningC, FaceA, FaceB, Closing, Simplicity}.lean — the
% ALM-A1 … ALM-A12 chain.  Fork A realizes a general mixed-sign H² curvature
% profile through a convex clean bicircle with symbolic levels 1 < a < c
% chosen per κ, culminating in hyperbolicMixedConverse_reparam_deg1:
% realization up to an orientation-preserving C¹ DEGREE-ONE circle
% reparametrization; the degree-one datum is consumed by
% chap:hyperbolic_converse to remove the reparametrization.

\section{Overview}

Fork A proves the genuinely-negative $H^2$ four-vertex converse for a
\emph{general} profile: a $\kappa$ satisfying
\cref{def:mixed_sign_hyperbolic_four_vertex} (continuous, $2\pi$-periodic,
escape velocity $>1$ at the maxima, minima \emph{arbitrarily negative}) is
realized, up to an orientation-preserving $C^1$ circle reparametrization, as
the geodesic curvature of a simple closed curve in $H^2$ ($\varepsilon=-1$).
Unlike the Euclidean and spherical forks the clean model is not a perturbed
circle but a \emph{convex clean bicircle} — arcs of constant curvature at two
symbolic levels $1<a<c$ chosen per $\kappa$ inside the four-vertex overlap gap
above $\max 1$ — and the whole chain runs in \emph{arc length}, on the engine
of \cref{chap:hyperbolic_arclength}.  The stages, one section per Lean
sub-module group:
\begin{enumerate}
\item \textbf{A1--A3 (Bicircle).}  Symbolic quarter residual $(G_1,G_2)$,
  monotonicity and sign brackets, nested-IVT anchor existence
  (\cref{thm:exists_bicircle_anchor}).
\item \textbf{A4 (Anchor).}  The closed-form anchor curve: closure,
  continuity, confinement, phase monotonicity, chord margin.
\item \textbf{A5 (Node).}  The node layout: period $\Lambda=L+w_1+w_2+t$,
  node map $g_{w,t}$ with quasi-periodicity $g(s+\Lambda)=g(s)+2\pi$
  (\cref{lem:node_map_add_period} — the conjugation that later makes the
  composite reparametrization degree one), the arc-length profile
  $\kappa_{\mathrm{arc}}=\kappa\circ h_1\circ g_{w,t}$.
\item \textbf{A6--A7 (Transport).}  Five-leg Grönwall transport, global
  confinement, residual continuity in the layout dofs.
\item \textbf{A8 (Turning A/B/C).}  Node-map inverse and coupling, field
  estimates, the continuous turning root $t=\tau(w)$.
\item \textbf{A9--A10 (Faces + Closing).}  Junction-chain columns, the
  face-sign theorem, the Poincaré--Miranda closing
  (\cref{thm:layout_closing}).
\item \textbf{A11--A12 (Simplicity + capstone).}  Chord non-vanishing, the
  window bridge into \cref{def:arclength_h2_curvature}, and the capstone
  \cref{thm:hyperbolic_mixed_converse_reparam_deg1} with
  $\Psi=(h_1\circ g_{w^*,t^*})\circ\chi$.
\end{enumerate}
A structural point shaping every statement below: the closing constants
$C_1,\varepsilon_0$ and the simplicity margin $\mu$ are quantified
\emph{ahead of} the reparametrization $h_1$; the capstone fixes
$\varepsilon=\min(\varepsilon_0,\mu/C_1)$ first and only then builds $h_1$,
breaking the reparam/$\varepsilon$ fixed point.

\section{The symbolic bicircle anchor (ALM-A1--A3)}

The bicircle quarter is a 2-arc composition: a level-$a$ model arc of length
$L/8$ from the start state $(ih,\pi)$, then a level-$c$ model arc of length
$L/8$.  With $W_2$ the quarter endpoint and $\varphi(L/4)$ the quarter phase,
the \emph{quarter residual} is $G_1=\operatorname{Im}W_2$,
$G_2=\varphi(L/4)-3\pi/2$; the quarter closes onto the Klein symmetry axes
exactly when $G_1=G_2=0$.  The free data are the start height $h$ and the
length $L$.

\begin{definition}\leanok
[the convex $h$-window]
  \label{def:bicircle_window}
  \lean{Gluck.Hyperbolic.bicircleWindow}
  For $a>1$ the \emph{$h$-window} is $\{h\mid 0<h<1,\ 2ah\le 1+h^2\}$, i.e.\
  $h\le r_a=(1-h^2)/(2(a-h))$ (the first-arc radius), equivalently
  $h\le a-\sqrt{a^2-1}$; on it every term of the $G_2$ monotone-difference
  factoring is nonnegative.  The accompanying \emph{$L$-bracket} is
  $\bar L(a,h)=4\pi r_a$ (\texttt{Gluck.Hyperbolic.bicircleBracket}): at
  $L=\bar L$ the first arc sweeps exactly $\pi/2$, so $G_2(h,\bar L)>0$ while
  $G_2(h,0)=-\pi/2$; moreover $\bar L<4\pi$, discharging the $L\le4\pi$
  hypothesis of the node layout.
\end{definition}

\begin{lemma}\leanok
[parametric IVT with continuous root selection]
  \label{lem:continuous_root_of_strict_mono}
  \lean{Gluck.Hyperbolic.continuous_root_of_strictMono}
  If on a parameter set $S$ the bracket endpoints $l\le u$ move continuously,
  $F\,x$ is continuous and strictly monotone on $[l\,x,u\,x]$ with locked
  endpoint signs $F\,x\,(l\,x)<0<F\,x\,(u\,x)$, and $F\,\cdot\,y$ is
  continuous on $S$ at each interior height $y$, then some $\rho$ continuous
  on $S$ selects the interior root: $\rho\,x\in(l\,x,u\,x)$,
  $F\,x\,(\rho\,x)=0$.
\end{lemma}

\begin{proof}
  \leanok
  IVT gives the root at each $x$; strict monotonicity makes it unique, and
  continuity of the selection follows by trapping $\rho$ between level sets
  of $F$ near each parameter, using continuity of $F\,\cdot\,y$ at interior
  heights.  Used twice: for the root $L^*(h)$ of $G_2$ here and for the
  turning root in \cref{lem:turning_root_continuous}.
\end{proof}

\begin{theorem}\leanok
[ALM-A3: symbolic anchor existence]
  \label{thm:exists_bicircle_anchor}
  \lean{Gluck.Hyperbolic.exists_bicircle_anchor}
  \uses{def:bicircle_window, lem:continuous_root_of_strict_mono}
  For every convex pair $1<a<c$ there are $h^*$ interior to the window, with
  $h^*\in[1/(10c),\,a-\sqrt{a^2-1}\,]$, and $L^*\in(0,\bar L(a,h^*))$ at
  which the quarter residual vanishes:
  $\operatorname{Im}W_2=0$ and $\varphi(L^*/4)=3\pi/2$.
\end{theorem}

\begin{proof}
  \uses{def:bicircle_window, lem:continuous_root_of_strict_mono}
  \leanok
  Nested IVT.  Inner: for each window $h$, $G_2(h,\cdot)$ is strictly
  increasing on $[0,\bar L]$ with $G_2(h,0)<0<G_2(h,\bar L)$, so
  \cref{lem:continuous_root_of_strict_mono} yields a \emph{continuous} root
  $L^*(h)$.  Outer: $h\mapsto G_1(h,L^*(h))$ is continuous on
  $[1/(10c),\,a-\sqrt{a^2-1}\,]$, negative at the left endpoint and positive
  at the right (at the boundary $r_a=h$ the $G_1>0$ endpoint sign fires
  exactly); an ordinary IVT gives $h^*$.
\end{proof}

\section{The anchor curve (ALM-A4)}

\begin{definition}\leanok
[the anchor curve]
  \label{def:anchor_curve}
  \lean{Gluck.Hyperbolic.anchorCurve}
  \uses{thm:exists_bicircle_anchor}
  The closed-form clean bicircle curve on $[0,L]$ at anchor data $(h,L)$: the
  quarter is the explicit two-arc constant-curvature composition; the half
  extends it by the conjugate Klein reflection
  $X(z,\varphi)=(\bar z,3\pi-\varphi)$; the full period by the central
  symmetry $\rho(z,\varphi)=(-z,\varphi+\pi)$.  The branch joints agree
  exactly because of the anchor equations of
  \cref{thm:exists_bicircle_anchor}, so the curve is closed, globally
  continuous, and has strictly increasing phase (file-private
  \texttt{anchorCurve\_continuous}, \texttt{anchorCurve\_phase\_strictMonoOn}).
\end{definition}

\begin{definition}\leanok
[the anchor confinement radius]
  \label{def:anchor_confine_radius}
  \lean{Gluck.Hyperbolic.anchorConfineRadius}
  \uses{def:anchor_curve}
  The explicit radius $R(a,c)=1-(a-1)(c-1)/(20c^2)<1$ confining the anchor
  curve strictly inside the Poincaré disk (file-private
  \texttt{anchorCurve\_confined}).  Two further A4 exports feed later stages:
  chord non-vanishing of every proper sub-arc, and the nonconstructive
  \emph{compact chord margin} — for each band width $\ell_0\in(0,L/2]$ a
  margin $m>0$ with $m(\tau-t)\le\|\int_t^\tau e^{i\varphi}\|$ whenever
  $\ell_0\le\tau-t\le L-\ell_0$ (compactness minimum of the nonvanishing
  chord; file-private \texttt{layout\_chord\_margin}).  The same template,
  with only the continuity input changed, gives the clean-layout chord margin
  consumed by \cref{lem:layout_chord_ne_zero}.
\end{definition}

\section{The node layout (ALM-A5)}

The anchor curve is a $(c,a,c,a,c)$ pattern of constant-curvature legs.  To
run a genuine profile on it, arc length is laid out by a \emph{node layout}:
interior widths $w_1,w_2$ and a terminal width $t$ deform the leg lengths,
giving the three degrees of freedom the closing needs ($2$ for $z$-closure
$+$ $1$ for turning).

\begin{definition}\leanok
[node density, breakpoints and period]
  \label{def:node_density}
  \lean{Gluck.Hyperbolic.nodeDensity}
  The \emph{node density} $\rho_{w,t}$ is a sum of $\Lambda$-periodic
  trapezoidal pulses (rescaled \texttt{clampTent}s of the Euclidean closing
  family) over the five layout legs, with breakpoints $s_1=\tfrac L8$,
  $s_2=\tfrac{3L}8+w_1$, $s_3=\tfrac{5L}8+w_1$, $s_4=\tfrac{7L}8+w_1+w_2$ and
  period $\Lambda=\texttt{nodePeriod}=L+w_1+w_2+t$
  (\texttt{Gluck.Hyperbolic.nodePeriod}).  On the layout box
  $|w_1|,|w_2|,|t|\le L/16$ it is continuous, strictly positive,
  $\Lambda$-periodic, with full-period integral exactly $2\pi$.
\end{definition}

\begin{definition}\leanok
[the node map $g_{w,t}$]
  \label{def:node_map}
  \lean{Gluck.Hyperbolic.nodeMap}
  \uses{def:node_density, def:integral_reparam}
  The running integral of the node density (an \texttt{integralReparam},
  \cref{def:integral_reparam}) anchored at $g(0)=3\pi/4$, the mid-$c$-arc
  point of the reference step.  On the box it is a $C^1$, strictly increasing
  reparametrization carrying the five layout legs onto the $\theta$-quarters
  of the canonical step
  $\kappa^*=\texttt{stepCurvature}\,c\,a\,0\,(\pi/2)\,\pi\,(3\pi/2)$
  (\cref{def:four_arc_step_function}).
\end{definition}

\begin{lemma}\leanok
[quasi-periodicity of the node map]
  \label{lem:node_map_add_period}
  \lean{Gluck.Hyperbolic.nodeMap_add_period}
  \uses{def:node_map, def:node_density}
  On the layout box ($0<L\le4\pi$, $|w_1|,|w_2|,|t|\le L/16$), for all $s$:
  \[
    g_{w,t}(s+\Lambda)\;=\;g_{w,t}(s)+2\pi .
  \]
  This conjugation law — period $\Lambda$ in arc length maps to period $2\pi$
  in step angle — is what makes the composite reparametrization of the
  capstone a \emph{degree-one} circle map: it composes with the shift laws of
  $h_1$ and of the linear window map $\chi$ in
  \cref{thm:hyperbolic_mixed_converse_reparam_deg1}.
\end{lemma}

\begin{proof}
  \uses{def:node_density}
  \leanok
  Split $\int_0^{s+\Lambda}\rho$ at $s$; the piece $\int_s^{s+\Lambda}\rho$
  equals the full-period integral by $\Lambda$-periodicity of the density,
  and the per-leg support integrals sum to exactly $2\pi$.
\end{proof}

\begin{definition}\leanok
[the arc-length curvature profile]
  \label{def:kappa_arc}
  \lean{Gluck.Hyperbolic.kappaArc}
  \uses{def:node_map}
  $\kappa_{\mathrm{arc}}=(\kappa\circ h_1)\circ g_{w,t}$: the four-vertex
  profile, pre-composed with the $L^1$-reparametrization $h_1$ (which makes
  $\kappa\circ h_1$ $L^1$-close to the clean step) and laid out in arc length
  by the node map.  Continuous, $\Lambda$-periodic
  (\cref{lem:node_map_add_period} composed with the shift law of $h_1$ and
  the periodicity of $\kappa$), bounded by any global bound of $\kappa$ —
  the profile the A6 true flow runs on.  Its clean counterpart is
  $\kappa^*\circ g_{w,t}$, the piecewise-constant $(c,a,c,a,c)$ profile.
\end{definition}

\begin{lemma}\leanok
[ALM-A5 capstone: comp-$L^1$ estimate]
  \label{lem:kappa_arc_comp_L1}
  \lean{Gluck.Hyperbolic.kappaArc_comp_L1}
  \uses{def:kappa_arc, def:node_map, def:four_arc_step_function,
        lem:node_map_add_period}
  On the layout box,
  \[
    \int_0^{\Lambda}\bigl|\kappa_{\mathrm{arc}}-\kappa^*\!\circ g_{w,t}\bigr|
    \;\le\;\frac L\pi\int_0^{2\pi}\bigl|\kappa\circ h_1-\kappa^*\bigr| .
  \]
\end{lemma}

\begin{proof}
  \uses{lem:node_map_add_period}
  \leanok
  Change of variables $\theta=g_{w,t}(s)$: the density is bounded below by
  $\pi/L$ on the box, so $ds\le(L/\pi)\,d\theta$, and one arc-length period
  maps onto one $\theta$-period by \cref{lem:node_map_add_period}.
\end{proof}

\section{Five-leg transport and residual continuity (ALM-A6--A7)}

\begin{definition}\leanok
[the true layout flow]
  \label{def:layout_flow}
  \lean{Gluck.Hyperbolic.layoutFlow}
  \uses{def:kappa_arc, def:anchor_curve, def:anchor_confine_radius}
  The \texttt{arcFlow} of $\kappa_{\mathrm{arc}}$ from the anchor mid-$c$
  start (state \texttt{layoutStart}, phase $5\pi/2$), at truncation radius
  $R'=\texttt{layoutConfineRadius}=(1+R_{\mathrm{clean}})/2$, horizon $2L$
  (covering every box period $\Lambda\le2L$ uniformly), curvature bound $M$.
  Its comparison object is the five-leg \texttt{layoutClean} curve — the
  anchor legs re-cut at the node breakpoints — confined in the disk of radius
  $R_{\mathrm{clean}}(a,c)=1-m_0\bigl((a-1)/(2(c+1))\bigr)^5<1$, the anchor
  margin $m_0$ of \cref{def:anchor_confine_radius} decayed by the per-leg
  escape ratio five times (\texttt{layoutClean\_confined}).
\end{definition}

\begin{lemma}\leanok
[ALM-A6: five-leg Grönwall transport]
  \label{lem:layout_trajectory_close}
  \lean{Gluck.Hyperbolic.layoutTrajectory_close}
  \uses{def:layout_flow, lem:kappa_arc_comp_L1}
  For anchor data on the window $\times$ bracket with the phase anchor
  equation, and continuous $2\pi$-periodic $\kappa$ with $|\kappa|\le M$:
  there is $C_1=C_1(a,c,L,M)>0$, uniform over the layout box and over $h_1$,
  with
  \[
    \bigl\|\Phi_{\mathrm{true}}(\sigma)-\Phi_{\mathrm{clean}}(\sigma)\bigr\|
    \;\le\;C_1\int_0^{2\pi}\bigl|\kappa\circ h_1-\kappa^*\bigr|
    \qquad(\sigma\in[0,\Lambda]).
  \]
\end{lemma}

\begin{proof}
  \uses{lem:kappa_arc_comp_L1}
  \leanok
  Chain the arc-flow Grönwall inequality across the five legs, each against
  the exact constant-level model-arc solution, with per-leg $L^1$ errors
  restricted from \cref{lem:kappa_arc_comp_L1}.  The constant —
  $5\,e^{5\,\mathrm{Lip}\,L}\cdot\frac2{1-R'^2}\cdot\frac L\pi$, with
  $\mathrm{Lip}$ the arc-field Lipschitz constant at radius $R'$, bound $M$
  — is explicit in the proof but exported existentially.
\end{proof}

\begin{lemma}\leanok
[ALM-A6: global confinement of the true flow]
  \label{lem:layout_flow_confined}
  \lean{Gluck.Hyperbolic.layoutFlow_confined}
  \uses{def:layout_flow, lem:layout_trajectory_close}
  If the true flow stays $b$-close to the clean curve on $[0,\Lambda]$ and
  $b\le(1-R_{\mathrm{clean}})/2$, then
  $\|z_{\mathrm{true}}\|\le R_{\mathrm{clean}}+b\le R'<1$ there; in
  particular the flow never reaches its truncation radius, so the clamped
  field equals the true field along the trajectory (the A12 window-bridge
  input).
\end{lemma}

\begin{proof}
  \leanok
  Triangle inequality from per-leg clean confinement plus the Grönwall gap
  $b$, and $R_{\mathrm{clean}}+b\le(1+R_{\mathrm{clean}})/2=R'$.
\end{proof}

\begin{definition}\leanok
[the layout closure residual]
  \label{def:layout_residual}
  \lean{Gluck.Hyperbolic.layoutResidual}
  \uses{def:layout_flow}
  $\mathcal R(w_1,w_2,t)
    =\Phi_{\mathrm{true}}(\Lambda)-\bigl(z(0),\varphi(0)+2\pi\bigr)$:
  endpoint state minus the one-full-turn closure target.  The first component
  is the $z$-closure residual (A10 consumes its real and imaginary parts);
  the second is the turning residual (the A8 root variable; on the anchor
  locus the turning target is $9\pi/2$).  $\mathcal R=0$ iff the true flow
  closes with total turning $2\pi$.
\end{definition}

\begin{lemma}\leanok
[ALM-A7: residual continuity on the layout box]
  \label{lem:layout_residual_continuous}
  \lean{Gluck.Hyperbolic.layoutResidual_continuousOn}
  \uses{def:layout_residual, lem:layout_trajectory_close}
  $(w_1,w_2,t)\mapsto\mathcal R(w_1,w_2,t)$ is jointly continuous on the
  layout box $|w_1|,|w_2|,|t|\le L/16$.
\end{lemma}

\begin{proof}
  \uses{lem:layout_trajectory_close}
  \leanok
  The endpoint state is continuous in the box dofs by the parametric
  Grönwall squeeze (the joint $(w,t)$-continuity ladder of the Transport
  module); the target is constant.
\end{proof}

\section{The turning nest (ALM-A8)}

The 3-dof closing is solved nested: first a continuous root $t=\tau(w)$ of
the turning residual, then a 2-D Poincaré--Miranda argument in
$w=(w_1,w_2)$.  Turning monotonicity needs the reparametrized profile
\emph{pointwise} close to $c$ on the terminal plateau — a strengthening of
the $L^1$ conclusion.

\begin{lemma}\leanok
[ALM-A8.0: plateau-pointwise $L^1$ reparametrization]
  \label{lem:bicircle_L1_reparam_pointwise}
  \lean{Gluck.Hyperbolic.exists_bicircle_L1_reparam_pointwise}
  \uses{lem:exists_preliminary_reparam, def:four_arc_step_function}
  For $\kappa$ continuous and $2\pi$-periodic (no positivity) with crossing
  data $\kappa(\theta_1,\ldots,\theta_4)=(a,c,a,c)$ at ordered points within
  one period, and every $\varepsilon>0$: there is an orientation-preserving
  circle reparametrization $h_1$ (strictly monotone, $C^1$ with continuous
  positive derivative witness, $h_1(\theta+2\pi)=h_1(\theta)+2\pi$) with
  \[
    \int_0^{2\pi}\bigl|\kappa\circ h_1-\kappa^*\bigr|<\varepsilon
    \quad\text{and}\quad
    \bigl|\kappa(h_1(\theta))-c\bigr|\le\varepsilon
    \ \text{on }[\tfrac\pi2,\tfrac{3\pi}4].
  \]
\end{lemma}

\begin{proof}
  \uses{lem:exists_preliminary_reparam}
  \leanok
  A re-run of the frozen Euclidean plateau construction
  (\cref{lem:exists_preliminary_reparam}, plateau variant), the
  half-race-width shift left-aligning the plateaus with the left-closed step
  quarters; the pointwise clause is the plateau value on the closed
  second-quarter window.  No positivity is needed: the construction uses
  continuity only, and the $L^1$ upgrade replaces the positive global bound
  by a two-sided compactness bound.
\end{proof}

\begin{lemma}\leanok
[ALM-A8: strict monotonicity of the turning residual in $t$]
  \label{lem:turning_residual_strict_mono}
  \lean{Gluck.Hyperbolic.turningResidual_strictMono_t}
  \uses{def:layout_residual, def:node_map, lem:bicircle_L1_reparam_pointwise,
        lem:layout_trajectory_close}
  For every anchor datum there is $\varepsilon_0>0$, uniform over the box,
  such that whenever $h_1$ satisfies the $L^1$ tolerance \emph{and} the
  plateau clause at level $\varepsilon\le\varepsilon_0$, the map
  $t\mapsto\mathcal R(w_1,w_2,t)_2$ is strictly increasing on
  $[-L/16,L/16]$, for every fixed box $(w_1,w_2)$.
\end{lemma}

\begin{proof}
  \uses{lem:bicircle_L1_reparam_pointwise, lem:layout_trajectory_close}
  \leanok
  Extending the terminal plateau inserts flow time at level $c$ with no
  downstream legs (terminal-dof locality of the node map).  The flows at
  $t<t'$ are coupled by the mass-matching coupling
  $\psi=g_{t'}^{-1}\circ g_t$ — fixing $s_4$, carrying $\Lambda_t$ to
  $\Lambda_{t'}$, $O(t'-t)$-close to the identity in value and derivative —
  and a second-order rectangle Grönwall estimate bounds the coupled
  difference by $C_2\,\varepsilon\,(t'-t)$, while the clean turning gain is
  $\ge 2(c-R_{\mathrm{clean}})(t'-t)$ (the plateau clause keeps the true
  terminal curvature within $\varepsilon$ of $c$).  Take
  $\varepsilon_0=\min(\varepsilon_1,(c-R_{\mathrm{clean}})/(C_2+1))$.
\end{proof}

\begin{lemma}\leanok
[ALM-A8: the continuous turning root $\tau(w)$]
  \label{lem:turning_root_continuous}
  \lean{Gluck.Hyperbolic.turningRoot_continuous}
  \uses{lem:turning_residual_strict_mono, lem:layout_residual_continuous,
        lem:continuous_root_of_strict_mono}
  There are a root-box radius $W_0\in(0,L/16]$ and $\varepsilon_0>0$ such
  that for every admissible $h_1$ at tolerance $\varepsilon\le\varepsilon_0$
  there is a continuous $\tau$ on $\{|w_1|,|w_2|\le W_0\}$ with
  $\tau(w)\in(-L/16,L/16)$ and $\mathcal R(w_1,w_2,\tau(w))_2=0$.
\end{lemma}

\begin{proof}
  \uses{lem:turning_residual_strict_mono, lem:layout_residual_continuous,
        lem:continuous_root_of_strict_mono}
  \leanok
  Strict monotonicity (\cref{lem:turning_residual_strict_mono}), a sign
  bracket at $t=\pm L/16$ (clean turning gain minus transported error), and
  joint continuity (\cref{lem:layout_residual_continuous}) feed
  \cref{lem:continuous_root_of_strict_mono}.
\end{proof}

\section{Face signs and the Poincaré--Miranda closing (ALM-A9--A10)}

\begin{definition}\leanok
[the clean $z$-closure residual]
  \label{def:a9_residual}
  \lean{Gluck.Hyperbolic.a9Residual}
  \uses{def:anchor_curve, def:layout_flow}
  The clean closure residual as an explicit, $\tau_{\mathrm{clean}}$-free map
  of the interior dofs $(w_1,w_2)$: the fixed-phase endpoint
  $z+r(1+ie^{i\psi})$ of the terminal $c$-leg through the fourth layout node,
  minus the layout start point.  It vanishes at the anchor $(0,0)$; its two
  exact derivative columns at $w=0$ are computed through the junction-chain
  columns of the FaceA module (the A9.0 circle-coordinate transfer matrices
  and the A9.1 reduced anchor data).
\end{definition}

\begin{theorem}\leanok
[ALM-A10: the true flow closes]
  \label{thm:layout_closing}
  \lean{Gluck.Hyperbolic.exists_layout_closing}
  \uses{def:a9_residual, lem:turning_root_continuous,
        lem:layout_trajectory_close, lem:layout_flow_confined,
        def:layout_residual, lem:layout_residual_continuous}
  For anchor data $(h,L)$ with both anchor equations, and continuous
  $2\pi$-periodic $\kappa$ with $|\kappa|\le M$: there are $C_1>0$ (the
  transport constant, exposed ahead of the threshold) and $\varepsilon_0>0$
  such that for every admissible plateau-pointwise $h_1$ at tolerance
  $\varepsilon\le\varepsilon_0$ there is a box point $(w_1,w_2,t)$ where the
  true flow \emph{closes with total turning $2\pi$}
  ($\mathcal R(w_1,w_2,t)=0$), bundled with the $C_1\varepsilon$-closeness to
  the clean curve and the confinement $\|z(\sigma)\|\le R'<1$ on the closed
  period window — the shapes A11 and A12 consume.
\end{theorem}

\begin{proof}
  \uses{def:a9_residual, lem:turning_root_continuous,
        lem:layout_trajectory_close, lem:layout_flow_confined,
        lem:layout_residual_continuous}
  \leanok
  The face-sign theorem (file-private \texttt{cleanClosure\_face\_signs},
  A9.3) supplies an invertible linear recombination $(A,B)$, $(A',B')$ of
  \cref{def:a9_residual} and a box cap $W_1\le L/16$ in the recombined dofs
  $u=w_1+w_2$, $v=w_1-w_2$: every radius $W\le W_1$ carries a face margin
  $m>0$ and phase tolerance $\eta>0$ such that near turning closure the first
  recombined component is $\ge m$ on the $u=W$ face and $\le-m$ on $u=-W$,
  likewise the second in $v$.  Intersect with the A8 root box
  ($W=\min(W_0,W_1)$) and slice along $t=\tau(w)$
  (\cref{lem:turning_root_continuous}): the turning residual vanishes, and
  the $z$-residual of the \emph{true} flow inherits the clean face signs with
  margin $m/2$, since \cref{lem:layout_trajectory_close} bounds the
  true$-$clean gap by $C_1\varepsilon$ and $\varepsilon_0$ is assembled from
  the A8 threshold and the quotas $C_1\varepsilon\le\eta$,
  $M_c\,C_1\varepsilon\le m/2$,
  $C_1\varepsilon\le(1-R_{\mathrm{clean}})/2$.  The Poincaré--Miranda
  rectangle engine yields $(u^*,v^*)$ where both recombined components
  vanish; invertibility ($AB'-BA'\ne0$) recovers $z$-closure, $\tau$ the
  turning closure, and \cref{lem:layout_flow_confined} the confinement.
\end{proof}

\section{Simplicity and the up-to-reparam capstone (ALM-A11--A12)}

\begin{lemma}\leanok
[ALM-A11: simplicity transport]
  \label{lem:layout_chord_ne_zero}
  \lean{Gluck.Hyperbolic.layout_chord_ne_zero}
  \uses{def:layout_flow, def:anchor_confine_radius, thm:layout_closing}
  There is a margin $\mu>0$, exported ahead of $C_1$ and $\varepsilon$, such
  that for the closed true flow of \cref{thm:layout_closing} with
  $C_1\varepsilon\le\mu$, every proper sub-arc chord is nonzero:
  $\int_p^q e^{i\varphi_{\mathrm{true}}(s)}\,ds\ne0$ for
  $0\le p<q<\Lambda$.
\end{lemma}

\begin{proof}
  \leanok
  Three regimes against the short scale $\ell_0=\pi/(3C_2)$, with
  $C_2=2(M+1)/(1-R'^2)$ the true phase-speed bound: \emph{short} arcs
  ($q-p\le\ell_0$; phase deviation $\le\pi/3$, midpoint projection of the
  chord positive — this tolerates the negative dips of $\kappa$);
  \emph{near-full} arcs (complement short, exact closure flips the chord);
  \emph{mid} arcs (the clean-layout chord margin $m_0$ — the A4 template of
  \cref{def:anchor_confine_radius} with the layout continuity input —
  transported at cost $2C_1\varepsilon\le2\mu$).  $\mu$ is exported first so
  A12 can fix $\varepsilon\le\mu/C_1$.
\end{proof}

\begin{lemma}\leanok
[ALM-A12: window-solution exposure]
  \label{lem:layout_window_solution}
  \lean{Gluck.Hyperbolic.layout_windowSolution_exposed}
  \uses{def:layout_flow, def:kappa_arc, thm:layout_closing,
        lem:layout_chord_ne_zero}
  From the closed, confined, chord-nonvanishing true flow: there are
  $z:\mathbb R\to\mathbb C$, $\varphi:\mathbb R\to\mathbb R$ with
  $z'=e^{i\varphi}$, $\varphi'$ the arc-length angular speed of
  $\kappa_{\mathrm{arc}}$, $\|z\|<1$, $z(\Lambda)=z(0)$,
  $\varphi(\Lambda)=\varphi(0)+2\pi$, and all proper sub-arc chords nonzero —
  the window-solution package at window $\Lambda$.
\end{lemma}

\begin{proof}
  \leanok
  The layout flow runs at horizon $2L$, the bridge consumes horizon
  $\Lambda\le2L$; the two arc flows agree on $[0,\Lambda]$ by ODE uniqueness
  (\texttt{arcFlow\_unique}), so the A10/A11 closure, confinement and chord
  data transfer verbatim.  Strict confinement is
  \cref{lem:layout_flow_confined} plus $R'<1$.
\end{proof}

\begin{lemma}\leanok
[ALM-A12: $\kappa_{\mathrm{arc}}$ is an $H^2$ arc-length curvature]
  \label{lem:layout_arclength_h2}
  \lean{Gluck.Hyperbolic.layout_arcLengthH2Curvature}
  \uses{lem:layout_window_solution, def:arclength_h2_curvature,
        def:kappa_arc}
  Under the hypotheses of \cref{lem:layout_window_solution},
  $\kappa_{\mathrm{arc}}=\kappa\circ h_1\circ g_{w,t}$ satisfies
  $\texttt{ArcLengthH2Curvature}$ (\cref{def:arclength_h2_curvature}) at
  window $\Lambda=\texttt{nodePeriod}$.
\end{lemma}

\begin{proof}
  \uses{lem:layout_window_solution}
  \leanok
  Package the window solution through the (public) arc-length window bridge
  \texttt{arcLengthH2Curvature\_of\_windowSolution} of
  \cref{chap:hyperbolic_arclength}.
\end{proof}

\begin{theorem}\leanok
[capstone: hyperbolic mixed converse, up to a degree-one reparametrization]
  \label{thm:hyperbolic_mixed_converse_reparam_deg1}
  \lean{Gluck.Hyperbolic.hyperbolicMixedConverse_reparam_deg1}
  \uses{def:mixed_sign_hyperbolic_four_vertex, def:spaceform_realizes,
        def:is_simple_closed, lem:exists_abab_levels,
        thm:exists_bicircle_anchor, lem:bicircle_L1_reparam_pointwise,
        thm:layout_closing, lem:layout_chord_ne_zero,
        lem:layout_arclength_h2, lem:node_map_add_period,
        thm:arclength_h2_converse, lem:arclength_h2_converse_at}
  Let $\kappa$ satisfy \cref{def:mixed_sign_hyperbolic_four_vertex} — in
  particular the minima may be \emph{arbitrarily negative}.  Then there are
  $z$ and $\Psi$ with
  \[
    \Psi\in C^1,\quad \Psi'>0,\quad
    \Psi(t+2\pi)=\Psi(t)+2\pi,\quad
    \texttt{IsSimpleClosed}\;z,\quad
    \texttt{Realizes}\,(-1)\;z\;(\kappa\circ\Psi).
  \]
  The third conjunct — $\Psi$ is a \emph{degree-one} circle map — is the
  datum \cref{chap:hyperbolic_converse} consumes
  (\cref{thm:realizes_of_reparam_degree_one}) to remove the
  reparametrization.  Since $H^2$ has no metric rescaling, the period is
  co-constructed rather than normalized (mirroring
  \cref{lem:realizes_h2_of_reparam}).
\end{theorem}

\begin{proof}
  \uses{lem:exists_abab_levels, thm:exists_bicircle_anchor,
        lem:bicircle_L1_reparam_pointwise, thm:layout_closing,
        lem:layout_chord_ne_zero, lem:layout_arclength_h2,
        lem:node_map_add_period, thm:arclength_h2_converse,
        lem:arclength_h2_converse_at}
  \leanok
  Split the hypothesis disjunction.  \emph{Constant branch}:
  $\kappa\equiv c>1$; the explicit escape-velocity hyperbolic circle
  (\texttt{hyperbolicCircle\_realizes}) with $\Psi=\mathrm{id}$.
  \emph{Four-vertex branch}: choose convex clean levels $1<a<b$ interior to
  the overlap gap above $\max 1$ (midpoints of the window halves around the
  hypothesis value $c$) and $(a,b,a,b)$ crossing data by
  \cref{lem:exists_abab_levels}; fix the symbolic anchor $(h^*,L^*)$ by
  \cref{thm:exists_bicircle_anchor} and a global bound $M$.  Quantify the
  reparam-uniform constants \emph{first}: the margin $\mu$
  (\cref{lem:layout_chord_ne_zero}) and $C_1,\varepsilon_0$
  (\cref{thm:layout_closing}); set
  $\varepsilon=\min(\varepsilon_0,\mu/C_1)$ — breaking the
  reparam/$\varepsilon$ fixed point.  Build the plateau-pointwise $h_1$ at
  tolerance $\varepsilon$ (\cref{lem:bicircle_L1_reparam_pointwise}), close
  the flow at some $(w_1^*,w_2^*,t^*)$ (\cref{thm:layout_closing}), transport
  simplicity (\cref{lem:layout_chord_ne_zero}, since
  $C_1\varepsilon\le\mu$), and certify
  $\texttt{ArcLengthH2Curvature}\,(\kappa\circ h_1\circ g_{w^*,t^*})$
  (\cref{lem:layout_arclength_h2}).  The arc-length converse
  (\cref{thm:arclength_h2_converse}, at-window form
  \cref{lem:arclength_h2_converse_at}) produces a simple closed curve
  realizing $\kappa\circ h_1\circ g_{w^*,t^*}\circ\chi$ with
  $\chi(u)=(\Lambda/2\pi)u$ the linear window reparametrization.  Set
  $\Psi=(h_1\circ g_{w^*,t^*})\circ\chi$: it is $C^1$ with positive
  derivative ($h_1$'s positive derivative witness, positive node density,
  $\chi'=\Lambda/2\pi>0$), and the three shift laws compose to degree one:
  $\chi(u+2\pi)=\chi(u)+\Lambda$, then $g(s+\Lambda)=g(s)+2\pi$
  (\cref{lem:node_map_add_period}), then
  $h_1(\theta+2\pi)=h_1(\theta)+2\pi$, giving
  $\Psi(u+2\pi)=\Psi(u)+2\pi$.
\end{proof}

\begin{corollary}\leanok
[up-to-reparam form, degree-one datum forgotten]
  \label{cor:hyperbolic_mixed_converse_family}
  \lean{Gluck.Hyperbolic.hyperbolicMixedConverse}
  \uses{thm:hyperbolic_mixed_converse_reparam_deg1}
  A $\texttt{MixedSignHyperbolicFourVertex}$ profile is realized, up to an
  orientation-preserving $C^1$ reparametrization, as the geodesic curvature
  of a simple closed curve in the hyperbolic plane.
\end{corollary}

\begin{proof}
  \uses{thm:hyperbolic_mixed_converse_reparam_deg1}
  \leanok
  Weakening of \cref{thm:hyperbolic_mixed_converse_reparam_deg1} forgetting
  the degree-one conjunct.
\end{proof}
